% GENERATED FILE. Edit canonical lessons, then run npm run generate:books.
% canonical-source-hash: fnv1a64:78d3ed9594c40a1f
% canonical-lessons: PA-C287-bund, PA-C287-bhijjna, PA-C287-harh, PA-C287-tapakna, PA-C287-chikkar, PA-R287-review-281-287

\chapter{When the Rain Comes}
\label{ch:pa-a2-when-the-rain-comes}
\hlchaptermodality{287}

\begin{chapteropening}
\textbf{By the end of this chapter:} \emph{I can say what the rain will bring.}

Complete the last lesson of chapter 287: I can say what the rain will bring.
\end{chapteropening}

\section[\texorpdfstring{\pa{ਬੂੰਦ} (būnd)}{ਬੂੰਦ (būnd)}]{\pa{ਬੂੰਦ} (būnd) — a drop}

\label{lesson:PA-C287-bund}



\begin{quote}
Before the new one: say the Punjabi for a squall, then the Punjabi for to flash.
\end{quote}

\subsection*{You'll want to know: \pa{ਬੂੰਦ}}
\textbf{\pa{ਬੂੰਦ}} — \emph{būnd} — \textquotedblleft{}a drop\textquotedblright{}.

\begin{quote}
\pa{ਛੱਤ} \pa{ਤੋਂ} \pa{ਪਾਣੀ} \pa{ਦੀ} \pa{ਇੱਕ} \pa{ਬੂੰਦ} \pa{ਡਿੱਗੀ।}
\end{quote}

\begin{quote}
\emph{chhatt ton pāṇī dī ikk būnd ḍiggī.}
\end{quote}

\begin{quote}
A drop of water fell from the roof.
\end{quote}

A feminine noun.

\begin{cousinweb}
\textbf{\pa{ਬੂੰਦ}} (\emph{būnd}) is inherited from Sanskrit \emph{bindu}, \textquotedblleft{}drop, spot\textquotedblright{}; CDIAL lists Punjabi \emph{bund}, \textquotedblleft{}drop\textquotedblright{}, under it (CDIAL 9240).
\end{cousinweb}

\begin{grammarlens}[title={What you've built}]
One more word for saying what comes next.
\end{grammarlens}

\subsection*{Your turn}
\begin{itemize}
\raggedright
  \item \emph{Say it:} \emph{būnd}
  \item \emph{Say it:} \emph{būnd}, once more
  \item \emph{Say it:} \emph{chhatt ton pāṇī dī ikk būnd ḍiggī.} — A drop of water fell from the roof.
  \item \emph{Recall:} say the Punjabi for a squall, then the Punjabi for to flash, then say \emph{būnd} again
\end{itemize}

\begin{tcolorbox}[breakable,colback=teal!4,colframe=teal!35!black,title={Before you move on}]
What does \pa{ਬੂੰਦ} mean? (\textquotedblleft{}A drop\textquotedblright{}.) Say it once more.
\end{tcolorbox}
\section[\texorpdfstring{\pa{ਭਿੱਜਣਾ} (bhijjṇā)}{ਭਿੱਜਣਾ (bhijjṇā)}]{\pa{ਭਿੱਜਣਾ} (bhijjṇā) — to get wet}

\label{lesson:PA-C287-bhijjna}



\begin{quote}
Before the new one: say the Punjabi for to flash, then the Punjabi for a drop.
\end{quote}

\subsection*{You'll want to know: \pa{ਭਿੱਜਣਾ}}
\textbf{\pa{ਭਿੱਜਣਾ}} — \emph{bhijjṇā} — \textquotedblleft{}to get wet\textquotedblright{}.

\begin{quote}
\pa{ਜੇ} \pa{ਮੀਂਹ} \pa{ਆਵੇਗਾ} \pa{ਤਾਂ} \pa{ਅਸੀਂ} \pa{ਭਿੱਜਾਂਗੇ।}
\end{quote}

\begin{quote}
\emph{je mīṁh āvegā tāṁ asīṁ bhijjāṅge.}
\end{quote}

\begin{quote}
If the rain comes, we will get wet.
\end{quote}

Past (I …): \pa{ਮੈਂ} \pa{ਮੀਂਹ} \pa{ਵਿੱਚ} \pa{ਭਿੱਜਿਆ}, \emph{maiṁ mīṁh vich bhijjiā}.

\begin{cousinweb}
\textbf{\pa{ਭਿੱਜਣਾ}} (\emph{bhijjṇā}) is inherited from a reconstructed \emph{*bhiyajyatē}, \textquotedblleft{}is smeared, gets wet\textquotedblright{}; CDIAL lists Punjabi \emph{bhijjṇā}, \textquotedblleft{}to get wet\textquotedblright{}, under it (CDIAL 9502).
\end{cousinweb}

\begin{grammarlens}[title={What you've built}]
One more word for saying what comes next.
\end{grammarlens}

\subsection*{Your turn}
\begin{itemize}
\raggedright
  \item \emph{Say it:} \emph{bhijjṇā}
  \item \emph{Say it:} \emph{bhijjṇā}, once more
  \item \emph{Say it:} \emph{je mīṁh āvegā tāṁ asīṁ bhijjāṅge.} — If the rain comes, we will get wet.
  \item \emph{Recall:} say the Punjabi for to flash, then the Punjabi for a drop, then say \emph{bhijjṇā} again
  \item \emph{Write it:} copy \textbf{\pa{ਭਿੱਜਣਾ}} once, checking the whole word
\end{itemize}

\begin{tcolorbox}[breakable,colback=teal!4,colframe=teal!35!black,title={Before you move on}]
What does \pa{ਭਿੱਜਣਾ} mean? (\textquotedblleft{}To get wet\textquotedblright{}.) Say it once more.
\end{tcolorbox}
\section[\texorpdfstring{\pa{ਹੜ੍ਹ} (haṛh)}{ਹੜ੍ਹ (haṛh)}]{\pa{ਹੜ੍ਹ} (haṛh) — a flood}

\label{lesson:PA-C287-harh}



\begin{quote}
Before the new one: say the Punjabi for a drop, then the Punjabi for to get wet.
\end{quote}

\subsection*{You'll want to know: \pa{ਹੜ੍ਹ}}
\textbf{\pa{ਹੜ੍ਹ}} — \emph{haṛh} — \textquotedblleft{}a flood\textquotedblright{}.

\begin{quote}
\pa{ਜੇ} \pa{ਮੀਂਹ} \pa{ਨਹੀਂ} \pa{ਰੁਕੇਗਾ} \pa{ਤਾਂ} \pa{ਹੜ੍ਹ} \pa{ਆਵੇਗਾ।}
\end{quote}

\begin{quote}
\emph{je mīṁh nahīṁ rukegā tāṁ haṛh āvegā.}
\end{quote}

\begin{quote}
If the rain does not stop, there will be a flood.
\end{quote}

A masculine noun.

\begin{cousinweb}
\textbf{\pa{ਹੜ੍ਹ}} (\emph{haṛh}) has no settled source: CDIAL gives no Punjabi form, and the Punjabi origin remains unknown in the sources used here.
\end{cousinweb}

\begin{grammarlens}[title={What you've built}]
One more word for saying what comes next.
\end{grammarlens}

\subsection*{Your turn}
\begin{itemize}
\raggedright
  \item \emph{Say it:} \emph{haṛh}
  \item \emph{Say it:} \emph{haṛh}, once more
  \item \emph{Say it:} \emph{je mīṁh nahīṁ rukegā tāṁ haṛh āvegā.} — If the rain does not stop, there will be a flood.
  \item \emph{Recall:} say the Punjabi for a drop, then the Punjabi for to get wet, then say \emph{haṛh} again
  \item \emph{Write it:} copy \textbf{\pa{ਹੜ੍ਹ}} once, checking the whole word
\end{itemize}

\begin{tcolorbox}[breakable,colback=teal!4,colframe=teal!35!black,title={Before you move on}]
What does \pa{ਹੜ੍ਹ} mean? (\textquotedblleft{}A flood\textquotedblright{}.) Say it once more.
\end{tcolorbox}
\section[\texorpdfstring{\pa{ਟਪਕਣਾ} (ṭapakṇā)}{ਟਪਕਣਾ (ṭapakṇā)}]{\pa{ਟਪਕਣਾ} (ṭapakṇā) — to drip; to leak}

\label{lesson:PA-C287-tapakna}



\begin{quote}
Before the new one: say the Punjabi for to get wet, then the Punjabi for a flood.
\end{quote}

\subsection*{You'll want to know: \pa{ਟਪਕਣਾ}}
\textbf{\pa{ਟਪਕਣਾ}} — \emph{ṭapakṇā} — \textquotedblleft{}to drip; to leak\textquotedblright{}.

\begin{quote}
\pa{ਮੀਂਹ} \pa{ਵਿੱਚ} \pa{ਪੁਰਾਣੀ} \pa{ਛੱਤ} \pa{ਟਪਕਦੀ} \pa{ਹੈ।}
\end{quote}

\begin{quote}
\emph{mīṁh vich purāṇī chhatt ṭapakdī hai.}
\end{quote}

\begin{quote}
The old roof drips in the rain.
\end{quote}

Past (it …): \pa{ਉਹ} \pa{ਛੱਤ} \pa{ਟਪਕੀ}, \emph{oh chhatt ṭapkī}.

\begin{cousinweb}
\textbf{\pa{ਟਪਕਣਾ}} (\emph{ṭapakṇā}) goes back to a reconstructed \emph{*ṭapp}, \textquotedblleft{}drip, drop\textquotedblright{}, an imitative root; CDIAL gives no Punjabi form, but lists Hindi \emph{ṭapaknā}, \textquotedblleft{}to drip\textquotedblright{} (CDIAL 5444).
\end{cousinweb}

\begin{grammarlens}[title={What you've built}]
One more word for saying what comes next.
\end{grammarlens}

\subsection*{Your turn}
\begin{itemize}
\raggedright
  \item \emph{Say it:} \emph{ṭapakṇā}
  \item \emph{Say it:} \emph{ṭapakṇā}, once more
  \item \emph{Say it:} \emph{mīṁh vich purāṇī chhatt ṭapakdī hai.} — The old roof drips in the rain.
  \item \emph{Recall:} say the Punjabi for to get wet, then the Punjabi for a flood, then say \emph{ṭapakṇā} again
  \item \emph{Write it:} copy \textbf{\pa{ਟਪਕਣਾ}} once, checking the whole word
\end{itemize}

\begin{tcolorbox}[breakable,colback=teal!4,colframe=teal!35!black,title={Before you move on}]
What does \pa{ਟਪਕਣਾ} mean? (\textquotedblleft{}To drip; to leak\textquotedblright{}.) Say it once more.
\end{tcolorbox}
\section[\texorpdfstring{\pa{ਚਿੱਕੜ} (chikkaṛ)}{ਚਿੱਕੜ (chikkaṛ)}]{\pa{ਚਿੱਕੜ} (chikkaṛ) — mud}

\label{lesson:PA-C287-chikkar}



\begin{quote}
Before the new one: say the Punjabi for a flood, then the Punjabi for to drip.
\end{quote}

\subsection*{You'll want to know: \pa{ਚਿੱਕੜ}}
\textbf{\pa{ਚਿੱਕੜ}} — \emph{chikkaṛ} — \textquotedblleft{}mud\textquotedblright{}.

\begin{quote}
\pa{ਮੀਂਹ} \pa{ਤੋਂ} \pa{ਬਾਅਦ} \pa{ਸੜਕ} \pa{ਉੱਤੇ} \pa{ਚਿੱਕੜ} \pa{ਹੋਵੇਗਾ।}
\end{quote}

\begin{quote}
\emph{mīṁh ton bāad saṛak utte chikkaṛ hovegā.}
\end{quote}

\begin{quote}
After the rain there will be mud on the road.
\end{quote}

A masculine noun.

\begin{cousinweb}
\textbf{\pa{ਚਿੱਕੜ}} (\emph{chikkaṛ}) goes back to Sanskrit \emph{cikka}, \textquotedblleft{}gummy matter\textquotedblright{}; CDIAL lists Punjabi \emph{cikkaṛ}, \textquotedblleft{}mud\textquotedblright{}, under it (CDIAL 4780).
\end{cousinweb}

\begin{grammarlens}[title={What you've built}]
One more word for saying what comes next.
\end{grammarlens}

\subsection*{Your turn}
\begin{itemize}
\raggedright
  \item \emph{Say it:} \emph{chikkaṛ}
  \item \emph{Say it:} \emph{chikkaṛ}, once more
  \item \emph{Say it:} \emph{mīṁh ton bāad saṛak utte chikkaṛ hovegā.} — After the rain there will be mud on the road.
  \item \emph{Recall:} say the Punjabi for a flood, then the Punjabi for to drip, then say \emph{chikkaṛ} again
  \item \emph{Write it:} copy \textbf{\pa{ਚਿੱਕੜ}} once, checking the whole word
\end{itemize}

\begin{tcolorbox}[breakable,colback=teal!4,colframe=teal!35!black,title={Before you move on}]
What does \pa{ਚਿੱਕੜ} mean? (\textquotedblleft{}Mud\textquotedblright{}.) Say it once more.
\end{tcolorbox}
\section[(dialogue)]{Looking back: from Plans for the Holidays to When the Rain Comes}

\label{lesson:PA-R287-review-281-287}



\begin{quote}
Say the Punjabi for a plan and an intention.
\end{quote}

\subsection*{Your turn}
Say these aloud:

\begin{itemize}
\raggedright
  \item a plan, an intention, advice, to put off, at once
  \item a suitcase, a bedding roll, a cloth shoulder bag, to keep safe, to come out
  \item a route, a companion, to arrive, to get stuck, to sway
  \item to get off, to take down, a push, last, to get to
  \item to return, to go in, to come to an end, to get tired, dawn
  \item a spell of steady rain, to thunder, to pour down, a squall, to flash
  \item a drop, to get wet, a flood, to drip, mud
\end{itemize}

\begin{tcolorbox}[breakable,colback=teal!4,colframe=teal!35!black,title={Before you move on}]
A plan? (\textbf{\pa{ਪ੍ਰੋਗਰਾਮ}}, \emph{progrām}.) An intention? (\textbf{\pa{ਇਰਾਦਾ}}, \emph{irādā}.) Advice? (\textbf{\pa{ਸਲਾਹ}}, \emph{salāh}.) To put off? (\textbf{\pa{ਟਾਲਣਾ}}, \emph{ṭālṇā}.) At once? (\textbf{\pa{ਤੁਰੰਤ}}, \emph{turant}.) A suitcase? (\textbf{\pa{ਅਟੈਚੀ}}, \emph{aṭaichī}.) A bedding roll? (\textbf{\pa{ਬਿਸਤਰਾ}}, \emph{bisatrā}.) A cloth shoulder bag? (\textbf{\pa{ਝੋਲਾ}}, \emph{jholā}.) To keep safe? (\textbf{\pa{ਸਾਂਭਣਾ}}, \emph{sāmbhṇā}.) To come out? (\textbf{\pa{ਨਿਕਲਣਾ}}, \emph{nikalṇā}.) A route? (\textbf{\pa{ਰਾਹ}}, \emph{rāh}.) A companion? (\textbf{\pa{ਸਾਥੀ}}, \emph{sāthī}.) To arrive? (\textbf{\pa{ਪਹੁੰਚਣਾ}}, \emph{pahunchṇā}.) To get stuck? (\textbf{\pa{ਫਸਣਾ}}, \emph{phasṇā}.) To sway? (\textbf{\pa{ਹਿੱਲਣਾ}}, \emph{hillṇā}.) To get off? (\textbf{\pa{ਉਤਰਨਾ}}, \emph{utarnā}.) To take down? (\textbf{\pa{ਉਤਾਰਨਾ}}, \emph{utārnā}.) A push? (\textbf{\pa{ਧੱਕਾ}}, \emph{dhakkā}.) Last? (\textbf{\pa{ਆਖ਼ਰੀ}}, \emph{ākhrī}.) To get to? (\textbf{\pa{ਅੱਪੜਨਾ}}, \emph{appaṛnā}.) To return? (\textbf{\pa{ਪਰਤਣਾ}}, \emph{paratṇā}.) To go in? (\textbf{\pa{ਵੜਨਾ}}, \emph{vaṛnā}.) To come to an end? (\textbf{\pa{ਮੁੱਕਣਾ}}, \emph{mukkṇā}.) To get tired? (\textbf{\pa{ਥੱਕਣਾ}}, \emph{thakkṇā}.) Dawn? (\textbf{\pa{ਪਹੁ}}, \emph{pahu}.) A spell of steady rain? (\textbf{\pa{ਝੜੀ}}, \emph{jhaṛī}.) To thunder? (\textbf{\pa{ਗੱਜਣਾ}}, \emph{gajjṇā}.) To pour down? (\textbf{\pa{ਵਰ੍ਹਨਾ}}, \emph{varhnā}.) A squall? (\textbf{\pa{ਝੱਖੜ}}, \emph{jhakkhaṛ}.) To flash? (\textbf{\pa{ਚਮਕਣਾ}}, \emph{chamakṇā}.) A drop? (\textbf{\pa{ਬੂੰਦ}}, \emph{būnd}.) To get wet? (\textbf{\pa{ਭਿੱਜਣਾ}}, \emph{bhijjṇā}.) A flood? (\textbf{\pa{ਹੜ੍ਹ}}, \emph{haṛh}.) To drip? (\textbf{\pa{ਟਪਕਣਾ}}, \emph{ṭapakṇā}.) Mud? (\textbf{\pa{ਚਿੱਕੜ}}, \emph{chikkaṛ}.)
\end{tcolorbox}
